\section{Obligations of each property}\label{app:obligations}

Table~\ref{tab:obligations} lists, per property of \S\ref{sec:properties},
what each component must do, what it rests on, and how it is established.

\begin{table*}[h]
  \caption{Obligations of each property and their status.}
  \label{tab:obligations}
  \small
  \begin{tabularx}{\textwidth}{@{}l>{\raggedright\arraybackslash}p{0.12\textwidth}>{\raggedright\arraybackslash}X>{\raggedright\arraybackslash}p{0.15\textwidth}>{\raggedright\arraybackslash}p{0.09\textwidth}>{\raggedright\arraybackslash}p{0.25\textwidth}@{}}
    \toprule
    Property & Circuit & Guardian and registry & Perch & Assumptions & Status \\
    \midrule
    1 Soundness & C1--C6 & checks 1--12; the constructor's circuit-id check; registry activation, grace, revocation and scope, with C6's hashing & statement built from controller state & A1a, A1b, A2, A3, A4, A5, A7 & proven over the model with named breaks and per-proof extraction; C1--C6 and the checks tested; parsing tested on real mail and fuzzing; registry hashing cross-checked \\
    2 Privacy & public inputs exactly $(\Cm, \digest_{hi}, \digest_{lo}, \mathit{root})$ & stores $\Cm$ and a salted pin root; emits nothing & --- & A3, A4, A6 & proven in the ideal model; structural lemma proven; the ideal model's step to the real primitives assumed \\
    3 Replay & C5 binds $\digest$ & checks 1, 2, 5, 6, 8, 12 & one vote per guardian per attempt and domain or per change digest; ids never reused; epochs grow & A2, A3, A4 & proven with the SHA-256 break named; the encoding facts tested \\
    4 One factor & --- & the guardian is one address & quorum, delay, cancellation & A9; A12 for the owner's veto & proven over all traces; wallet policy proven as a predicate and tested in code \\
    \bottomrule
  \end{tabularx}
\end{table*}

\section{Sender authentication per provider}\label{app:providers}

A1b is a property of each provider's submission path. Table~\ref{tab:providers}
summarizes a review of public documentation and published measurements, read,
not tested; the repository's review
(\nolinkurl{research/email-recovery/eval/parser/sender-authentication.md})
cites a source for each cell\claim{provider-authenticates-sender}. Where
\texttt{h=} names \texttt{from} more often than it occurs (``oversigned''), a
second \texttt{From} lands in the signed block and the circuit refuses it, and
the circuit refuses a \texttt{From} naming several addresses whatever the
provider does (\S\ref{sec:design-circuit}). The measurements are from 2020 and
2021, and the providers responded to them; a controlled submission test from
test accounts would settle column~1 today, and needs real mail sent through
real providers.

\begin{table*}[h]
  \caption{The three conditions of A1b per launch provider: (1) two
    \texttt{From} fields or several addresses at submission; (2) sending as
    another address of the provider's domain; (3) reassignment of deleted or
    inactive addresses, which $\code$ covers.}
  \label{tab:providers}
  \small
  \begin{tabularx}{\textwidth}{@{}>{\raggedright\arraybackslash}p{0.13\textwidth}>{\raggedright\arraybackslash}p{0.09\textwidth}>{\raggedright\arraybackslash}X>{\raggedright\arraybackslash}p{0.21\textwidth}>{\raggedright\arraybackslash}p{0.2\textwidth}@{}}
    \toprule
    Provider & \texttt{h=} & (1) & (2) & (3) \\
    \midrule
    Gmail & oversigned & 2020: checked the angle-address the circuit binds~\cite{chen-composition-kills}; today unknown & needs a confirmation link & not reassigned \\
    Outlook, Hotmail & once & 2020: the one strict path of eight tested~\cite{chen-composition-kills}; variants bypassed checks~\cite{shen-weak-links}; today unknown & another user's mailbox: unknown & not reassigned \\
    iCloud & once & 2020: checked only the first of several \texttt{From} fields~\cite{chen-composition-kills}; today unknown & another user's mailbox: unknown & unknown \\
    Yahoo & oversigned & 2020: checked only the first address~\cite{chen-composition-kills} and accepted a second field~\cite{shen-weak-links}; the circuit refuses both forms & needs a verification message & inactive IDs released in 2013; today unknown \\
    Proton & oversigned & not tested by either study; unknown & addresses are created on the account & not reassigned \\
    Fastmail & oversigned & 2020: accepted arbitrary \texttt{From} addresses~\cite{chen-composition-kills}; today unknown & needs a code sent to the address & reassigned; covered by $\code$ \\
    GMX, WEB.DE & oversigned & not tested by either study; unknown & another user's mailbox: unknown & reassigned after a year of inactivity; covered by $\code$ \\
    \bottomrule
  \end{tabularx}
\end{table*}

\section{Open questions}\label{app:open}

Each open decision, with its current setting and the alternatives.
\begin{enumerate}
  \item \emph{Pin semantics.} Additive, fixed at construction, salted,
    expiring. A restrictive pin, or requiring pin and registry, makes a
    planted key useless but stops recovery at each provider rotation until a
    new guardian is deployed; a pin the owner can change lets whoever holds
    the owner key and $\code$ pin a key of their own.
  \item \emph{Pin lifetime.} 30 days of close time, capped at 365: past the
    7-day timelock with margin, and a short window in which a revoked key
    stays usable through a pin. The cap keeps the enrolled key, and a revoked
    one, usable longest; about 10 days narrows the window further.
  \item \emph{Default allowed actions.} Every action. Initiation and
    cancellation only, or initiation only, with wallet rule 4 preventing a
    lock-out; under a guardian-only quorum of one an inbox thief with the
    code can also approve cancellations of the owner's attempts, which are not
    alerted.
  \item \emph{Where $\code$ is kept.} The encrypted copy in the Sent folder
    and the kit. Plain text in the enrollment email recovers best but lets
    the provider compute $\Cm$; a kit alone keeps no copy at the provider but
    is lost with the kit.
  \item \emph{Registry parameters.} A 24-hour grace, capped at 7 days (7
    days strands fewer proofs but keeps a blanked key acceptable for a week);
    an emergency revocation threshold of 2 of 5 (under the full threshold two
    members can neither block a live key nor stop a planted one); the
    members, attesters and their resolvers are not named.
  \item \emph{Provider details.} Each provider's signing domains, its signed
    block for long addresses and display names, and whether its Sent copy
    keeps the signature, against real outbound mail.
  \item \emph{The recovery address.} The address, the production mail
    exchanger in the enclave with TLS terminating there, and the domain's
    MTA-STS policy.
  \item \emph{Alert code derivation.} Keyed from $\code$, so any device
    holding it, a phisher's included, shows it. A key from the passkey's PRF
    output lets the owner's device authenticate alerts but not a recovering
    one; a code from the public digest needs no key but lets the sending
    service and the provider link an alert to its attempt.
  \item \emph{On-chain enrollment proof.} Only the wallet checks the
    enrollment email; a constructor that verifies a proof of it would make
    $\Cm$ checkable by third parties, at the cost of a second circuit.
  \item \emph{Guardian deployment.} Directly from the wallet's submitter,
    checked by code hash and configuration; or through a factory that fixes
    them on chain; and who pays on mainnet.
  \item \emph{The push channel's state.} Checked on this device at once and
    through the vault with the passkey, as fresh as the weekly probe; reading
    it on every open costs a passkey prompt.
  \item \emph{Custom domains.} Larger address and domain bounds change the
    commitment's packing and the leaf format.
  \item \emph{Custodians, the enclave and its image registry.} Who the
    custodians are (one outside the operator), their production protocol,
    the enclave's production binding, and the image registry's signers.
  \item \emph{Outages longer than RPC retention.} Reconciliation and a lag
    alarm are built; a monitor outside the watcher is not, nor a choice of
    whether the email-plus-delay floor should exceed the longest tolerated
    outage.
  \item \emph{Accounts without an email guardian} have no $\code$ to key
    their alert codes; a key in the recovery kit is one alternative.
  \item \emph{What the alert host learns} (which accounts have contacts):
    whether to blur it with cover traffic or batched reads.
  \item \emph{Failover duplicates.} The second sending service has no
    idempotency key, so an accepted request whose answer was lost can be
    delivered twice.
  \item \emph{Probe visibility.} Each weekly probe shows a quiet
    notification, because WebKit revokes subscriptions whose pushes show
    nothing; how long subscriptions survive on each platform is not measured.
  \item \emph{A lowered target close time.} Existing email-plus-delay members
    then last less than 7 days until rebuilt; the wallet does not yet notice
    and ask owners to rebuild, or keep a margin above the floor.
  \item \emph{Sender authentication} (Table~\ref{tab:providers}): every
    provider's handling of two \texttt{From} fields today, sending as another
    mailbox at three providers, and reassignment at two.
  \item \emph{Keeping contracts live.} Who funds the keeper on mainnet,
    whether the wallet extends an account's own entries, and the registry's
    tree paths.
\end{enumerate}
